(1) 令 $x_{\rm st}=\Re[Xe^{i\omega t}]$，代入得 $X=F_0/[k-m\omega^2+ib\omega]$。故幅度 $A=F_0/\sqrt{(k-m\omega^2)^2+b^2\omega^2}$，相对驱动滞后角 $\delta=\arg(k-m\omega^2+ib\omega)$，满足 $\tan\delta=b\omega/(k-m\omega^2)$，需按象限选正确相位。
(2) 最大位移幅度对应分母平方 $D=(k-m\omega^2)^2+b^2\omega^2$ 最小。对 $\omega>0$ 求导得 $dD/d\omega=2\omega[-2m(k-m\omega^2)+b^2]$，所以 $\omega_{\rm disp}^2=k/m-b^2/(2m^2)$。它低于无阻尼本征频率 $\omega_0=\sqrt{k/m}$；当阻尼足够大时此正频率共振峰不再存在。
(3) 稳态速度为 $\dot x=-A\omega\sin(\omega t-\delta)$，故阻尼平均功率 $\overline P_b=b\langle\dot x^2\rangle=b\omega^2A^2/2$。外力平均功率 $\overline P_{\rm in}=\langle F_0\cos\omega t\,\dot x\rangle=F_0A\omega\sin\delta/2$；由 $\sin\delta=b\omega A/F_0$ 可知两者严格相等。
(4) 弱阻尼邻近 $\omega_0$，有 $k-m\omega^2\simeq-2m\omega_0(\omega-\omega_0)$，功率谱近似洛伦兹峰。半高发生在 $|\omega-\omega_0|\simeq b/(2m)$，全宽 $\Delta\omega\simeq b/m$。品质因子 $Q\simeq\omega_0/\Delta\omega=m\omega_0/b$，要求高 $Q$ 才使用该近似。
